\chapter*{Słowniczek}
\addcontentsline{toc}{chapter}{Słowniczek}
\markboth{Słowniczek}{}
Ten słowniczek objaśnia podstawowe pojęcia i skróty z podręcznika. Kody ról w zespole (KM, KS, OT i inne) i kody formularzy (F0--F14, ICD) opisuje dodatek~\ref{dod:formularze}.

\renewcommand{\arraystretch}{1.0}
\small
\section*{Pojęcia}
\noindent\begin{tabularx}{\linewidth}{@{}>{\bfseries}P{33mm}L@{}}
\toprule
Balon stratosferyczny & Zestaw z powłoką wypełnioną gazem nośnym, olinowaniem, spadochronem i gondolą. Wznosi się do stratosfery (zwykle 30--\SI{35}{km}), potem powłoka pęka albo odcinacz kończy lot.\\
Powłoka & Worek z lateksu (albo z folii), który napełniasz gazem nośnym. Powłoka lateksowa rośnie w czasie wznoszenia i pęka na pułapie.\\
Gaz nośny & Gaz lżejszy od powietrza, który unosi balon: hel albo wodór.\\
Siła nośna & Różnica między ciężarem powietrza wypartego przez powłokę a ciężarem gazu i powłoki. Podajemy ją w gramach albo kilogramach.\\
Siła nośna szyjki & (\emph{neck lift}) Siła, z którą powłoka ciągnie szyjkę w górę: masa wszystkiego, co wisi pod powłoką, plus nadwyżka, która daje prędkość wznoszenia. Mierzysz ją obciążnikiem.\\
Obciążnik & Ciężar o masie równej sile nośnej szyjki. W czasie tankowania wisi na szyjce. Gdy balon go unosi, tankowanie jest zakończone.\\
Szyjka & Wąska końcówka powłoki. Przez nią tankujesz gaz. Po tankowaniu zamykasz ją i mocujesz do niej olinowanie.\\
Tankowanie & Napełnianie powłoki gazem nośnym z butli przez reduktor.\\
Reduktor & Zawór na butli. Obniża ciśnienie gazu i ustawia przepływ.\\
Olinowanie & Linki i złącza, które łączą powłokę, spadochron i gondolę.\\
Krętlik & Obrotowe złącze w olinowaniu. Nie pozwala skręcić się linkom.\\
Słabe ogniwo & Element olinowania między powłoką a spadochronem, który zrywa się przy ustalonej sile.\\
Odcinacz & Urządzenie, które odcina powłokę od reszty zestawu i kończy lot, np. po ustalonym czasie.\\
Gondola & Lekka, izolowana skrzynka z elektroniką i ładunkiem.\\
Ładunek & Eksperyment albo urządzenie, które leci na platformie.\\
Platforma & Gondola z elektroniką, olinowanie, spadochron i odcinacz. Platforma daje ładunkom lot, zasilanie i łączność. W tym podręczniku: Stardust NEXT.\\
Zestaw lotny & Wszystko, co leci: powłoka, olinowanie, spadochron, gondola i ładunki.\\
Segment lotny, naziemny & Zestaw lotny oraz wszystko, co zostaje na ziemi: stanowisko startowe, stacje odbiorcze, samochody i zespół.\\
Pułap & Największa wysokość lotu.\\
Pęknięcie & (\emph{burst}) Moment, w którym powłoka pęka. Zestaw zaczyna opadać na spadochronie.\\
Dryf & (\emph{float}) Lot na prawie stałej wysokości, gdy balon przestaje się wznosić.\\
Predykcja trajektorii & Obliczenie trasy lotu i miejsca lądowania z prognozy wiatru.\\
Tracker & Nadajnik w gondoli. Wysyła pozycję z GPS i telemetrię.\\
Telemetria & Dane z lotu, które nadajnik wysyła przez radio w czasie lotu.\\
Stacja odbiorcza & Odbiornik z anteną. Odbiera sygnał trackera i przekazuje dane do sieci śledzenia.\\
Digipeater, iGate & Stacje sieci APRS. Digipeater nadaje ponownie odebrane pakiety. iGate przekazuje pakiety do internetu (APRS-IS).\\
Znak wywoławczy & Identyfikator stacji krótkofalarskiej, np. SP4EVA. SSID to numer po myślniku, np. SP4EVA-11.\\
Go / no-go & Decyzja „leci / nie leci” przed startem.\\
Pościg & Jazda samochodami za balonem do miejsca lądowania.\\
\bottomrule
\end{tabularx}

\clearpage
\section*{Skróty}
\noindent\begin{tabularx}{\linewidth}{@{}>{\bfseries}P{20mm}L@{}}
\toprule
AIS & Służba informacji lotniczej (\emph{Aeronautical Information Service}).\\
AMC & Komórka zarządzania przestrzenią powietrzną (\emph{Airspace Management Cell}).\\
APRS & Radiowy system pakietowy do przesyłania pozycji (\emph{Automatic Packet Reporting System}). APRS-IS -- jego sieć w internecie.\\
ASM1 & Komórka PAŻP, która przyjmuje zgłoszenia sondowań atmosfery.\\
ATZ & Strefa ruchu lotniska (\emph{Aerodrome Traffic Zone}).\\
AUP, UUP & Plan użytkowania przestrzeni powietrznej i jego aktualizacja (\emph{Airspace Use Plan}, \emph{Updated AUP}).\\
BEXUS, REXUS & Programy lotów balonowych i rakietowych dla studentów (\emph{Balloon / Rocket Experiments for University Students}).\\
CDR, PDR & Przeglądy projektu: krytyczny i wstępny (\emph{Critical / Preliminary Design Review}).\\
CTR, MCTR & Strefa kontrolowana lotniska cywilnego i wojskowego (\emph{Control Zone}).\\
ESA & Europejska Agencja Kosmiczna.\\
FRR & Przegląd gotowości do lotu (\emph{Flight Readiness Review}).\\
GPS & Globalny system nawigacji satelitarnej.\\
ICD & Dokument kontroli interfejsu platforma--ładunek (\emph{Interface Control Document}).\\
ISM & Pasma radiowe dla urządzeń przemysłowych, naukowych i medycznych. Nie wymagają licencji, ale mają limity mocy.\\
KML & Format pliku z trasą (np. Google Earth).\\
LDPE & Polietylen o niskiej gęstości -- folia na powłoki eksperymentalne.\\
LNA & Wzmacniacz o niskim szumie (\emph{Low-Noise Amplifier}).\\
LoRa, LoRaWAN & Modulacja radiowa o dużym zasięgu i małej mocy oraz sieć bramek, która jej używa.\\
NOTAM & Komunikat dla pilotów, np. o sondowaniu atmosfery (\emph{Notice to Airmen}).\\
NOPATH, WIDE2-1 & Ścieżki pakietów APRS. NOPATH: pakiet nie jest powtarzany. WIDE2-1: digipeatery powtarzają pakiet.\\
OBC & Komputer pokładowy (\emph{On-Board Computer}).\\
OC & Ubezpieczenie od odpowiedzialności cywilnej.\\
OGN & Sieć śledzenia szybowców (\emph{Open Glider Network}).\\
P, R, D & Strefy przestrzeni powietrznej: zakazana, ograniczona, niebezpieczna.\\
PAŻP & Polska Agencja Żeglugi Powietrznej.\\
PCB & Płytka drukowana.\\
PMR446 & Radiotelefony na pasmo \SI{446}{MHz}, bez licencji.\\
RTTY & Prosta telemetria tekstowa przez radio (\emph{Radio Teletype}).\\
SDR, RTL-SDR & Radio programowe (\emph{Software Defined Radio}). RTL-SDR: tani odbiornik SDR na USB.\\
SED & Dokumentacja eksperymentu w programie REXUS/BEXUS (\emph{Student Experiment Documentation}).\\
SERA & Wspólne europejskie zasady ruchu lotniczego, rozporządzenie (UE) nr 923/2012.\\
SSDV, SSTV & Przesyłanie zdjęć przez radio: w pakietach cyfrowych (SSDV) albo analogowo (SSTV).\\
SWR & Współczynnik fali stojącej -- miara dopasowania anteny.\\
T0, T$-$x & Chwila startu oraz czas przed startem, np. T$-$2 dni.\\
UTC & Uniwersalny czas koordynowany. Czas lokalny w Polsce: UTC+1 zimą, UTC+2 latem.\\
VHF, UHF & Zakresy radiowe 30--\SI{300}{MHz} i 300--\SI{3000}{MHz}. W nich są pasma krótkofalarskie \SI{2}{m} (\SI{144}{MHz}) i \SI{70}{cm} (\SI{430}{MHz}).\\
\bottomrule
\end{tabularx}
\normalsize
